\subsection{Pilot đa phương thức trên dữ liệu công khai}\label{sec:multimodal-pilot}
Phép bổ sung này dùng các cặp ảnh--tweet của CrisisMMD có nhãn hai modality thống nhất, với nhiệm vụ nhị phân informative/not informative~\cite{alam2018crisismmd,ofli2020multimodal}. Đây là thông tin hữu ích cho ứng phó thiên tai, khác nhãn khẩn cấp và lợi ích cứu hộ. Nguồn mô tả corpus tiếng Anh; chưa có đánh giá tiếng Việt. Pilot được thực hiện riêng trên CPU, không thay dữ liệu hoặc kết quả RQ1--RQ3.

Protocol khóa trước huấn luyện và dự đoán test. Giữ sự kiện lũ Sri Lanka làm test, bão Maria làm validation và năm sự kiện còn lại làm train. Ghép ba split gốc rồi chia theo sự kiện, nên kết quả không trực tiếp đối sánh với official split hoặc số liệu trong bài tham khảo. Chọn theo hash, tối đa 3.000 train và 600 validation, giữ test đủ điều kiện; mỗi nhóm tweet/text chỉ chọn một cặp. Nối thêm nhóm theo hash ảnh và dHash có khoảng cách Hamming không quá 4; loại 48 mẫu thuộc nhóm xuyên split, không lấy bù. Còn \MMTrainN{} train, \MMDevN{} validation và \MMTestN{} test trong \MMTestGroups{} nhóm bản trùng. Test có 182 informative/519 not informative, khác train 2.188/783; không điều chỉnh ngưỡng theo prior test.

Nhánh text học TF-IDF từ câu tự do, unigram/bigram, tối đa 20.000 chiều; vocabulary chỉ fit train. Nhánh ảnh lấy 576 đặc trưng từ MobileNetV3 Small ImageNet cố định, chưa fine-tuning. Hai block chuẩn hóa L2, ghép và chia $\sqrt{2}$; một LR regularized được huấn luyện trên nhãn chung. Mô hình kết hợp có trọng số được học nhưng chỉ là fusion tuyến tính cộng các đặc trưng, chưa học tương tác hoặc fine-tuning chung hai encoder. Ba classifier ảnh/text/ghép cùng thử $C\in\{0{,}1;1;10\}$, class weight cân bằng và ngưỡng $0{,}5$; chọn macro-F1 validation, hòa chọn $C$ nhỏ hơn. Lớp đa số được chọn từ train; trung bình xác suất hai nhánh là đối chứng kết hợp cố định, chưa hiệu chuẩn xác suất.

Bảng~\ref{tab:multimodal-pilot} trình bày kết quả trên cùng test IDs; recall $(+)$ là recall informative. Hai hàng cuối che một block của fusion khi suy luận, khác với huấn luyện lại nhánh đơn ở hai hàng text/ảnh.
% Generated by scripts/export_multimodal.py; do not edit.
\begin{table}[htbp]\centering
\caption{Pilot CrisisMMD trên sự kiện lũ Sri Lanka; hai hàng cuối che modality khi suy luận}\label{tab:multimodal-pilot}
\begin{tabular}{lrrr}\toprule
Cấu hình & Macro-F1 (\%) & Accuracy (\%) & Recall (+) (\%) \\ \midrule
Lớp đa số & 20,61 & 25,96 & 100,00 \\
Text TF-IDF + LR & 76,60 & 79,46 & 85,71 \\
Ảnh frozen + LR & 76,19 & 79,32 & 82,97 \\
Trung bình xác suất & 77,12 & 80,17 & 84,07 \\
Nối đặc trưng + LR & 78,35 & 81,17 & 86,81 \\
Fusion thiếu ảnh & 48,55 & 75,46 & 6,04 \\
Fusion thiếu text & 72,77 & 76,32 & 77,47 \\
\bottomrule\end{tabular}
\end{table}


Fusion đạt macro-F1 \MMFusionFOne\%, so với text \MMTextFOne\%, ảnh \MMImageFOne\% và trung bình xác suất \MMMeanFOne\%. Khi che nhánh ảnh của fusion cố định, recall informative chỉ còn 6,04\%, trong khi nhánh text được huấn luyện riêng đạt 85,71\%. Mất một block tạo đầu vào ngoài điều kiện train; kết quả không hỗ trợ khả năng chịu thiếu modality của fusion. CLI dùng classifier đơn tương ứng khi thiếu nguồn, nhưng chính sách chuyển nhánh này chưa được đánh giá trong vận hành.

Bảng~\ref{tab:multimodal-effects} báo hiệu macro-F1 fusion trừ từng đối chứng. Bootstrap dùng cùng 2.000 draw theo \MMTestGroups{} nhóm để giữ ghép cặp. CI 95\% chưa hiệu chỉnh đa so sánh và chỉ có điều kiện trong một sự kiện; số nhóm không phải số sự kiện độc lập. Không diễn giải việc một CI loại trừ 0 thành ưu thế trên toàn bộ thiên tai hoặc mức ý nghĩa toàn họ.
% Generated by scripts/export_multimodal.py; do not edit.
\begin{table}[htbp]\centering
\caption{Hiệu macro-F1 của fusion nối đặc trưng: bootstrap ghép cặp theo nhóm bản trùng, CI 95\% chưa hiệu chỉnh}\label{tab:multimodal-effects}
\begin{tabular}{lrrr}\toprule
Đối chứng & Hiệu (đpt) & Cận dưới & Cận trên \\ \midrule
Text & 1,75 & -2,06 & 5,46 \\
Ảnh & 2,16 & -0,41 & 4,80 \\
Trung bình xác suất & 1,23 & -1,20 & 3,85 \\
\bottomrule\end{tabular}
\end{table}


Hiệu fusion so với text là \MMFusionTextDelta{} điểm phần trăm, CI $[\MMFusionTextLow;\MMFusionTextHigh]$; so với ảnh là \MMFusionImageDelta{} điểm phần trăm, CI $[\MMFusionImageLow;\MMFusionImageHigh]$. Cả ba CI đều chứa 0. Điểm ước lượng cao hơn trong pilot chưa đủ để kết luận fusion vượt trội hoặc hai cấu hình tương đương.

Đợt bổ sung cung cấp bằng chứng đã huấn luyện nhánh đọc văn bản và fusion đặc trưng, cùng ablation có target/test chung trong một benchmark công khai. Việc chỉ dùng cặp nhãn thống nhất tạo thiên lệch chọn mẫu; nguồn có lỗi mã hóa văn bản và chưa được kiểm ngôn ngữ từng hàng. Chưa có lặp nhiều sự kiện test hoặc nhiều seed học sâu, chưa đánh giá tiếng Việt, khẩn cấp, điều phối hoặc Android. CLI nghiên cứu local gắn đúng task informative; xác suất này không được ánh xạ vào $E_i$ hay xếp hạng cứu hộ. Nhánh cứu hộ tiếng Việt vẫn cần dữ liệu ghép cặp, nhãn độc lập và xác nhận chuyên môn.

\paragraph{Chuyển từ kết quả pilot sang thiết kế tiếng Việt.} Pilot trả lời được câu hỏi kỹ thuật rằng ba nhánh có thể được học và đối chiếu trên cùng nhãn informative và cùng \MMTestN{} ID. Nó cũng ghi nhận giới hạn khi một block vắng mặt, nên thiết kế mở rộng không được mặc định fusion luôn tốt hơn hoặc chịu thiếu nguồn. Thiết kế ở Mục~\ref{sec:vi-ablation-plan} dùng nhiệm vụ mới là có/không có yêu cầu hỗ trợ cứu hộ trong tin, với nhãn dự thảo và corpus chưa có. Các metric của bảng pilot không được sao sang bảng kế hoạch tiếng Việt; chỉ phương thức ghép ID, khóa test và đối chiếu cấu hình được dùng làm kinh nghiệm thiết kế.
